\section{Introduction}
\IEEEPARstart{E}{vent-related} potentials (ERPs) are time-locked measurements of evoked neural dynamics, and decoding them with machine learning has become a standard engineering route to brain-computer interfaces, clinical markers, and affective state estimation \cite{blankertz2011erp,roy2019review,craik2019review}. Every decoding pipeline places a scaling step between the recorded epoch and the model. A widespread choice scales each epoch and channel by its own statistics, subtracting the epoch mean and dividing by the epoch standard deviation. The step is motivated by numerical conditioning and by the large between-recording differences in EEG amplitude, and a recent normalization benchmark reports that window-level, within-channel normalization gives the highest supervised performance among the strategies it compared \cite{truong2025normalization}.

Preprocessing choices shape what a decoder measures. High-pass filtering displaces multivariate patterns into activity-silent periods \cite{vandriel2021highpass} and creates artifactual ERP effects of opposite polarity before a true effect \cite{tanner2015highpass}. A multiverse analysis of the ERP CORE dataset reports significant decoding clusters that extend into the baseline period for some forking paths \cite{kessler2025preprocessing}. Benchmarking studies in this journal show that EEG results are sensitive to preprocessing \cite{robbins2020sensitive} and that preprocessing depth changes deep-learning performance across six tasks \cite{delpup2025more}. Leakage through data partitioning is now widely documented \cite{brookshire2024leakage,kapoor2023leakage}. Normalization that uses each epoch's own statistics is a different kind of leakage. It moves information along the time axis inside a single observation, so subject-grouped partitioning and train-only scaling leave it untouched. Its closest precedent is the detrending analysis of \cite{vandriel2021highpass}, which showed empirically that a polynomial fitted to a whole trial displaces decodable information into the baseline and that masking the response during the fit prevents it. Subtracting an epoch mean is the zero-order case of such a fit. Three things remain open after that work, and they are the subject of this one: a closed form for what the normalization step writes into the pre-stimulus interval, the role of the dividing statistic, which no masking of a subtracted trend addresses, and a law for how large the effect is in a given decoding regime. The step also now sits by default in deep-learning toolchains. A recent foundation-model benchmark compares z-scoring, unit rescaling, and division by each trial's 95th amplitude percentile as its normalization choices \cite{wu2025adabrain}, whereas causal alternatives such as exponential moving standardization \cite{schirrmeister2017deep} were designed for online use.

This work isolates that mechanism and quantifies it. It grows out of the temporal encoding stage of ARSPI-Net \cite{lane2023arspinet,lane2024arspinet,lane2026dissertation}, a fixed leaky integrate-and-fire (LIF) reservoir \cite{maass2002lsm} whose perturbation properties on affective ERPs were characterized at IEEE BIBM 2026 \cite{lane2026bibm}. The spike-count window of that reservoir lies mostly before stimulus onset, and its input epochs were z-scored over their full length. A fixed recurrent system read out before the stimulus provides an unusually clean probe. Any condition information in its state at onset can only enter through the scaling step.

The paper makes four contributions. First, it proves two exact identities. Through the subtracted statistic, per-epoch z-scoring maps the pre-stimulus mean of a baseline-corrected epoch onto a sign-flipped, rescaled copy of the post-stimulus mean. Through the dividing statistic, it writes post-stimulus energy into every pre-stimulus variance, including under scale-only normalization. A causal encoder therefore receives post-stimulus information before the stimulus (Sec.~\ref{sec:theory}). Second, it derives an averaging law that predicts the size of the effect from single-trial statistics. Discriminability of the leaked feature grows as the square root of the number of averaged trials, which places trial-averaged, participant-level decoding (the regime of many clinical and individual-differences studies) in the high-leakage regime and single-trial decoding in the low-leakage regime. Third, it tests these predictions on 211 participants of an affective-picture ERP cohort and on five paradigms of the public ERP CORE dataset \cite{kappenman2021erpcore}, with linear readouts, the ARSPI-Net reservoir, and a trained EEGNet \cite{lawhern2018eegnet}. Fourth, it re-evaluates the BIBM 2026 reservoir result under scaling that uses no post-stimulus sample and gives a reporting rule with a pre-stimulus null test that any pipeline can run.
